% GENERATED FILE. Edit canonical lessons, then run npm run generate:books.
% canonical-source-hash: fnv1a64:9fb032eaee969c5f
% canonical-lessons: HI-C84-uska, HI-C84-hamara, HI-C84-tumhara, HI-C84-apna, HI-C84-mere-paas, HI-C84-mere-liye, HI-C84-repaso-kiska, HI-C84-sintesis-mere-paas

\chapter{Whose, And What A Postposition Does To It}
\label{ch:hi-possessive}
\hlchaptermodality{92}

\begin{chapteropening}
\textbf{By the end of this chapter:} \emph{I can say whose something is for any owner, use apna when the owner is the subject, and say what I have, using the oblique form a postposition demands.}

Ask what somebody has and answer, noticing that no word in either line means have and that every possessive has bent to -e because a postposition follows it.
\end{chapteropening}

\section[uskā]{uskā — one word for his, her and its}

\label{lesson:HI-C84-uska}



\begin{quote}
You have had \emph{merā}, \emph{merī}, \emph{mere} since the second chapter. Say all three. You have been able to say whose something is for exactly one person: your own.
\end{quote}

\subsection*{You'll want to know: uskā}
\emph{uskā} — \textquotedblleft{}his,\textquotedblright{} \textquotedblleft{}her,\textquotedblright{} or \textquotedblleft{}its.\textquotedblright{}

Hindi does not choose between those three, and the reason is the point of the lesson: \textbf{the possessive agrees with the thing owned, not with the owner.}

\begin{grammarlens}[title={Grammar Lens: English agrees with the owner, Hindi with the thing}]
This is the single biggest difference between the two systems, and it runs the opposite way from what an English speaker expects:

\noindent
\begin{tabularx}{\linewidth}{@{}>{\raggedright\arraybackslash}X>{\raggedright\arraybackslash}X@{}}
\toprule
\textbf{} & \textbf{what decides the form} \\
\midrule
English \emph{his} / \emph{her} & the \textbf{owner}'s gender \\
Hindi \emph{uskā} / \emph{uskī} & the \textbf{owned thing}'s gender \\
\bottomrule
\end{tabularx}

So \emph{uskā} before a masculine noun and \emph{uskī} before a feminine one — and whether the owner is a man or a woman changes nothing at all. A sentence with \emph{uskī kitāb} tells you the book is feminine and tells you \textbf{nothing} about whose it is.

This is not a new rule. It is the \textbf{-\hi{आ}} agreement you already have, doing its ordinary job: \emph{uskā} ends in \textbf{-\hi{आ}}, so it bends like every other \textbf{-\hi{आ}} word in the book.
\end{grammarlens}

\subsection*{Your turn}
Say these aloud:

\begin{itemize}
\raggedright
  \item \emph{uskā} — his, her, or its
  \item the form before a feminine noun (\emph{uskī})
  \item \emph{merā} and \emph{uskā} one after the other, and hear the same ending
  \item what \emph{uskī} tells you about the owner (nothing)
\end{itemize}

\begin{tcolorbox}[breakable,colback=teal!4,colframe=teal!35!black,title={Before you move on}]
Does \emph{uskā} mean \emph{his} or \emph{her}? (Either — Hindi does not distinguish.) What decides between \emph{uskā} and \emph{uskī}? (The gender of the thing owned.)
\end{tcolorbox}
\section[hamārā]{\hi{हमारा} — our, and the places Hindi prefers it to \textquotedblleft{}my\textquotedblright{}}

\label{lesson:HI-C84-hamara}



\begin{quote}
Say \emph{merā} and \emph{uskā}. Both end the same way and bend the same way, so the next word brings no new grammar — only a new owner.
\end{quote}

\subsection*{You'll want to know: \hi{हमारा}}
\textbf{\hi{हमारा}} (\emph{hamārā}) — \textquotedblleft{}our.\textquotedblright{}

Built on \textbf{\hi{हम}} (\emph{ham}), \textquotedblleft{}we,\textquotedblright{} which you have met. Before a feminine noun it is \emph{hamārī}; before a plural or an oblique it is \emph{hamāre}, exactly as \emph{merā} behaves.

\begin{grammarlens}[title={Grammar Lens: where Hindi says \textquotedblleft{}our\textquotedblright{} and English says \textquotedblleft{}my\textquotedblright{}}]
The grammar is easy; the usage is the thing worth carrying:

\noindent
\begin{tabularx}{\linewidth}{@{}>{\raggedright\arraybackslash}X>{\raggedright\arraybackslash}X>{\raggedright\arraybackslash}X@{}}
\toprule
\textbf{talking about} & \textbf{English habit} & \textbf{Hindi habit} \\
\midrule
the house you grew up in & my house & \textbf{hamārā} ghar \\
your family & my family & \textbf{hamārā} parivār \\
your country & my country & \textbf{hamārā} deś \\
\bottomrule
\end{tabularx}

A Hindi speaker describing their family will very often say \emph{hamārā parivār} where an English speaker says \emph{my family}, because the thing is understood as belonging to the household rather than to the speaker. Neither is wrong in either language — but the preference is real, and a listening paper will play you the Hindi one.
\end{grammarlens}

\subsection*{Your turn}
Say these aloud:

\begin{itemize}
\raggedright
  \item \emph{hamārā} — our
  \item the form before a feminine noun (\emph{hamārī})
  \item the word for \emph{we} that it is built on (\emph{ham})
  \item \emph{hamārā parivār}, and notice that an English speaker would have said \emph{my}
\end{itemize}

\begin{tcolorbox}[breakable,colback=teal!4,colframe=teal!35!black,title={Before you move on}]
What pronoun is \emph{hamārā} built on? (\emph{Ham}, we.) Which of \emph{uskā} and \emph{hamārā} tells you something about the owner? (\emph{Hamārā} — it says the owner is a group including the speaker.)
\end{tcolorbox}
\section[tumhārā]{\hi{तुम्हारा} — your, at the middle level of three}

\label{lesson:HI-C84-tumhara}



\begin{quote}
Say \emph{hamārā}. You have had \emph{āpkā} for \emph{your} since early on. This lesson gives you the other one, and says when each is right.
\end{quote}

\subsection*{You'll want to know: \hi{तुम्हारा}}
\textbf{\hi{तुम्हारा}} (\emph{tumhārā}) — \textquotedblleft{}your,\textquotedblright{} said to someone you are familiar with.

Built on \textbf{\hi{तुम}} (\emph{tum}), the middle of Hindi's three \emph{you}s.

\begin{grammarlens}[title={Grammar Lens: three levels, and the possessive that goes with each}]
\noindent
\begin{tabularx}{\linewidth}{@{}>{\raggedright\arraybackslash}X>{\raggedright\arraybackslash}X>{\raggedright\arraybackslash}X@{}}
\toprule
\textbf{pronoun} & \textbf{possessive} & \textbf{who} \\
\midrule
\hi{आप} \emph{āp} & \emph{āpkā} & a stranger, an elder, anyone respected \\
\hi{तुम} \emph{tum} & \textbf{tumhārā} & a friend, a sibling, a younger colleague \\
\hi{तू} \emph{tū} & \emph{terā} & intimate, or rude — rarely the right choice for a learner \\
\bottomrule
\end{tabularx}

The safe default for a learner is \textbf{\hi{आप}}, and the reason is asymmetric: using \emph{āp} with a close friend sounds a little formal, while using \emph{tum} with a stranger or an elder sounds \textbf{rude}. One error is a small stiffness; the other is an offence.

But \emph{tumhārā} is not optional to know. A listening paper is full of friends talking to each other, and it is the form they use.
\end{grammarlens}

\subsection*{Your turn}
Say these aloud:

\begin{itemize}
\raggedright
  \item \emph{tumhārā} — your, to a friend
  \item the form before a feminine noun (\emph{tumhārī})
  \item the pronoun it is built on (\emph{tum})
  \item which of the two \emph{your}s you would use with somebody's grandmother (\emph{āpkā})
\end{itemize}

\begin{tcolorbox}[breakable,colback=teal!4,colframe=teal!35!black,title={Before you move on}]
Which pronoun is \emph{tumhārā} built on? (\emph{Tum}.) Which of the two mistakes costs more — \emph{āp} to a friend, or \emph{tum} to an elder? (\emph{Tum} to an elder; the other is only a little formal.)
\end{tcolorbox}
\section[apnā]{\hi{अपना} — the word that takes over when the owner is the subject}

\label{lesson:HI-C84-apna}



\begin{quote}
Say \emph{merā} and \emph{tumhārā}. Now try to say \textquotedblleft{}I do my work\textquotedblright{} in Hindi using \emph{merā}, and read on to find out why that sentence is wrong.
\end{quote}

\subsection*{You'll want to know: \hi{अपना}}
\textbf{\hi{अपना}} (\emph{apnā}) — \textquotedblleft{}one's own.\textquotedblright{}

When the owner of a thing is \textbf{the subject of the same sentence}, Hindi replaces the ordinary possessive with \emph{apnā}. This is a requirement rather than a preference:

\noindent
\begin{tabularx}{\linewidth}{@{}>{\raggedright\arraybackslash}X>{\raggedright\arraybackslash}X@{}}
\toprule
\textbf{} & \textbf{} \\
\midrule
wrong & \emph{maiṅ \textbf{merā} kām kartā hūṅ} \\
right & \emph{maiṅ \textbf{apnā} kām kartā hūṅ} \\
\bottomrule
\end{tabularx}

Both would be \textquotedblleft{}I do my work.\textquotedblright{} The first is what an English speaker builds, and a Hindi ear hears it as a mistake straight away.

\begin{grammarlens}[title={Grammar Lens: one word covering every person}]
The saving is large, because \emph{apnā} does not care who the subject is:

\noindent
\begin{tabularx}{\linewidth}{@{}>{\raggedright\arraybackslash}X>{\raggedright\arraybackslash}X>{\raggedright\arraybackslash}X@{}}
\toprule
\textbf{subject} & \textbf{English} & \textbf{Hindi} \\
\midrule
I & my own & \textbf{apnā} \\
you & your own & \textbf{apnā} \\
he or she & their own & \textbf{apnā} \\
we & our own & \textbf{apnā} \\
\bottomrule
\end{tabularx}

Four English possessives, one Hindi word. It still agrees with the thing owned in the ordinary way — \emph{apnā kām}, \emph{apnī kitāb}, \emph{apne dost} — because it ends in \textbf{-\hi{आ}} like all the others.

And it removes a real ambiguity. \emph{Vah \textbf{uskī} kitāb paṛhtā hai} means he reads \textbf{somebody else's} book; \emph{vah \textbf{apnī} kitāb paṛhtā hai} means he reads \textbf{his own}. English \emph{he reads his book} cannot tell those apart without more words.
\end{grammarlens}

\subsection*{Your turn}
Say these aloud:

\begin{itemize}
\raggedright
  \item \emph{apnā} — one's own
  \item \emph{maiṅ apnā kām kartā hūṅ}, and stop yourself from reaching for \emph{merā}
  \item the difference between \emph{apnī kitāb} and \emph{uskī kitāb}
  \item the form before a plural or an oblique (\emph{apne})
\end{itemize}

\begin{tcolorbox}[breakable,colback=teal!4,colframe=teal!35!black,title={Before you move on}]
When does \emph{apnā} replace the ordinary possessive? (When the owner is the subject of the sentence.) Whose book is \emph{uskī kitāb} when the subject is \emph{he}? (Somebody else's.)
\end{tcolorbox}
\section[mere pās]{\hi{मेरे} \hi{पास} — Hindi has no word for \textquotedblleft{}have\textquotedblright{}}

\label{lesson:HI-C84-mere-paas}



\begin{quote}
Say \emph{merā}, \emph{merī}, \emph{mere}. You have had all three since the second chapter. This lesson shows you that the third one has a second job nobody mentioned.
\end{quote}

\subsection*{You'll want to know: \hi{मेरे} \hi{पास}}
\textbf{\hi{मेरे} \hi{पास}} (\emph{mere pās}) — \textquotedblleft{}I have.\textquotedblright{} Literally: \textbf{near me}.

\textbf{\hi{पास}} (\emph{pās}) means \textquotedblleft{}near, at hand.\textquotedblright{} Hindi has no verb meaning \emph{to have}, so possession is built out of location:

\begin{quote}
\hi{मेरे} \hi{पास} \hi{एक} \hi{किताब} \hi{है।} \emph{mere pās ek kitāb hai} — \textquotedblleft{}Near me there is a book\textquotedblright{} = \textbf{I have a book.}
\end{quote}

The thing owned is the \textbf{subject} of that sentence, which is why \emph{hai} agrees with \emph{kitāb} and not with the speaker.

\begin{grammarlens}[title={Grammar Lens: mere is not the plural here}]
This is the part that repays attention. You learned \emph{mere} as the \textbf{plural} of \emph{merā}. It is also the \textbf{oblique} — the form a word takes in front of a postposition — and the two happen to look identical:

\noindent
\begin{tabularx}{\linewidth}{@{}>{\raggedright\arraybackslash}X>{\raggedright\arraybackslash}X@{}}
\toprule
\textbf{phrase} & \textbf{why \emph{mere}} \\
\midrule
\emph{mere dost} & plural — my friends \\
\emph{mere pās} & oblique — \emph{pās} is a postposition \\
\bottomrule
\end{tabularx}

So \emph{mere pās} is not \textquotedblleft{}my nears.\textquotedblright{} It is \emph{merā} bending because something stands after it. Every \textbf{-\hi{आ}} possessive does this: \emph{uskā} $\to$ \emph{uske pās}, \emph{hamārā} $\to$ \emph{hamāre pās}, \emph{apnā} $\to$ \emph{apne pās}.

That single habit — \textbf{-\hi{आ}} becomes \textbf{-\hi{ए}} in front of a postposition — is worth more than the phrase itself, because it governs every postposition in the language.
\end{grammarlens}

\subsection*{Your turn}
Say these aloud:

\begin{itemize}
\raggedright
  \item \emph{mere pās ek kitāb hai} — I have a book
  \item the same about him or her (\emph{uske pās})
  \item the same about us (\emph{hamāre pās})
  \item the same about the subject's own (\emph{apne pās})
  \item what \emph{mere} is doing in \emph{mere pās} — plural, or oblique? (oblique)
\end{itemize}

\begin{tcolorbox}[breakable,colback=teal!4,colframe=teal!35!black,title={Before you move on}]
What does \emph{mere pās} literally say? (Near me.) What agrees with \emph{hai} in \emph{mere pās ek kitāb hai}? (\emph{Kitāb} — the thing owned is the subject.) What happens to an \textbf{-\hi{आ}} possessive before a postposition? (It becomes \textbf{-\hi{ए}}.)
\end{tcolorbox}
\section[mere liye]{\hi{मेरे} \hi{लिए} — the same bend, in front of a different postposition}

\label{lesson:HI-C84-mere-liye}



\begin{quote}
Say \emph{mere pās}. Say why it is \emph{mere} and not \emph{merā}. The word in this lesson works the same way, which is the whole point of putting them together.
\end{quote}

\subsection*{You'll want to know: \hi{मेरे} \hi{लिए}}
\textbf{\hi{मेरे} \hi{लिए}} (\emph{mere liye}) — \textquotedblleft{}for me.\textquotedblright{} \emph{mere sāth} — \textquotedblleft{}with me.\textquotedblright{}

Three phrases, one shape. The possessive goes oblique and the postposition follows:

\noindent
\begin{tabularx}{\linewidth}{@{}>{\raggedright\arraybackslash}X>{\raggedright\arraybackslash}X@{}}
\toprule
\textbf{phrase} & \textbf{meaning} \\
\midrule
mere pās & near me — \emph{I have} \\
mere liye & for me \\
mere \emph{sāth} & with me \\
\bottomrule
\end{tabularx}

\begin{grammarlens}[title={Grammar Lens: the pattern generalises to every person}]
Once the bend is automatic, the whole grid follows without memorising a single new form:

\noindent
\begin{tabularx}{\linewidth}{@{}>{\raggedright\arraybackslash}X>{\raggedright\arraybackslash}X>{\raggedright\arraybackslash}X@{}}
\toprule
\textbf{owner} & \textbf{oblique} & \textbf{with \emph{liye}} \\
\midrule
I & mere & mere liye \\
he, she & uske & uske liye \\
we & hamāre & hamāre liye \\
the subject's own & apne & apne liye \\
\bottomrule
\end{tabularx}

The last row is the one both practice papers reach for. Phrases like \emph{apne dost ke sāth} — \textquotedblleft{}with one's own friend\textquotedblright{} — put three of this chapter's ideas in four words: the reflexive possessive, the oblique bend, and a two-word postposition.
\end{grammarlens}

\subsection*{Your turn}
Say these aloud:

\begin{itemize}
\raggedright
  \item \emph{mere liye} — for me
  \item \emph{mere sāth} — with me
  \item the same two about him or her (\emph{uske liye}, \emph{uske sāth})
  \item \emph{apne dost ke sāth}, and name the three things happening in it
\end{itemize}

\begin{tcolorbox}[breakable,colback=teal!4,colframe=teal!35!black,title={Before you move on}]
Why is it \emph{mere liye} and not \emph{merā liye}? (A postposition follows, so the possessive goes oblique.) What is the oblique of \emph{apnā}? (\emph{Apne}.)
\end{tcolorbox}
\section[pā̃ch mālik, ek niyam]{pā̃ch mālik, ek niyam — five owners, and one bend that covers them all}

\label{lesson:HI-C84-repaso-kiska}



\begin{quote}
Name as many possessives as you can before reading the table. Aim for five.
\end{quote}

\begin{grammarlens}[title={Grammar Lens: the whole grid}]
Every one of these ends in \textbf{-\hi{आ}}, so every one bends the same way:

\noindent
\begin{tabularx}{\linewidth}{@{}>{\raggedright\arraybackslash}X>{\raggedright\arraybackslash}X>{\raggedright\arraybackslash}X@{}}
\toprule
\textbf{owner} & \textbf{masculine} & \textbf{feminine} \\
\midrule
my & merā & merī \\
your (respectful) & āpkā & āpkī \\
your (familiar) & tumhārā & tumhārī \\
his, her, its & uskā & uskī \\
our & hamārā & hamārī \\
the subject's own & apnā & apnī \\
\bottomrule
\end{tabularx}

Six forms in the first column and not one of them is irregular. The agreement is with the \textbf{thing owned} throughout, which is why nothing in the middle column tells you anything about the owner.
\end{grammarlens}

\begin{grammarlens}[title={Grammar Lens: the bend, and the two jobs of -\hi{ए}}]
Put a postposition after any of them and the \textbf{-\hi{आ}} becomes \textbf{-\hi{ए}}:

\noindent
\begin{tabularx}{\linewidth}{@{}>{\raggedright\arraybackslash}X>{\raggedright\arraybackslash}X@{}}
\toprule
\textbf{direct} & \textbf{before a postposition} \\
\midrule
merā & \textbf{mere} pās \\
uskā & \textbf{uske} liye \\
hamārā & \textbf{hamāre} ghar meṅ \\
apnā & \textbf{apne} dost ke \emph{sāth} \\
\bottomrule
\end{tabularx}

The form \textbf{-\hi{ए}} therefore has two jobs, and the sentence tells you which:

\begin{itemize}
\raggedright
  \item \emph{mere dost} — no postposition follows, so this is the \textbf{plural}: my friends.
  \item \emph{mere pās} — a postposition follows, so this is the \textbf{oblique}: near me.
\end{itemize}

Nothing in the spelling distinguishes them. What comes next does.
\end{grammarlens}

\subsection*{Your turn}
Say these aloud:

\begin{itemize}
\raggedright
  \item all six possessives in the masculine, then all six in the feminine
  \item the one that replaces the others when the owner is the subject (\emph{apnā})
  \item \emph{mere pās}, \emph{uske pās}, \emph{hamāre pās} in a run
  \item which job \textbf{-\hi{ए}} is doing in \emph{mere dost}, and which in \emph{mere pās}
\end{itemize}

\begin{tcolorbox}[breakable,colback=teal!4,colframe=teal!35!black,title={Before you move on}]
What does a possessive agree with — the owner or the thing owned? (The thing owned.) What are the two jobs of the \textbf{-\hi{ए}} ending? (Plural, and oblique before a postposition.)
\end{tcolorbox}
\section[āpke pās kyā hai?]{\hi{आपके} \hi{पास} \hi{क्या} \hi{है}? — asking what somebody has, without a verb for having}

\label{lesson:HI-C84-sintesis-mere-paas}



\begin{quote}
Say \emph{mere pās ek kitāb hai}. Now turn it into a question addressed to somebody you would call \emph{āp}, and notice that \emph{āpkā} will have to bend.
\end{quote}

\begin{grammarlens}[title={Grammar Lens: the question bends too}]
\noindent
\begin{tabularx}{\linewidth}{@{}>{\raggedright\arraybackslash}X>{\raggedright\arraybackslash}X@{}}
\toprule
\textbf{} & \textbf{sentence} \\
\midrule
question & \hi{आपके} \hi{पास} \hi{क्या} \hi{है}? \\
answer & \hi{मेरे} \hi{पास} \hi{एक} \hi{किताब} \hi{है।} \\
\bottomrule
\end{tabularx}

\emph{Āpkā} becomes \emph{āpke} for the same reason \emph{merā} becomes \emph{mere}: \textbf{\hi{पास}} stands after it. The question and the answer are the identical construction with the owner swapped.

There is no verb for \emph{have} anywhere in either line. \textbf{\hi{है}} is the ordinary copula, and it is agreeing with \emph{kitāb}.
\end{grammarlens}

\begin{tcolorbox}[breakable,skin=enhanced,colback=orange!6,colframe=orange!45!black,boxrule=0.5pt,arc=1mm,left=6pt,right=6pt,top=4pt,bottom=4pt,fonttitle=\bfseries,title={Read this}]
Read aloud. Every word has been taught:

\begin{quote}
— \hi{आपके} \hi{पास} \hi{क्या} \hi{है}? — \hi{मेरे} \hi{पास} \hi{एक} \hi{किताब} \hi{है।} — \hi{तुम्हारे} \hi{पास} \hi{क्या} \hi{है}? — \hi{मेरे} \hi{पास} \hi{एक} \hi{कलम} \hi{है।} — \hi{हमारे} \hi{घर} \hi{में} \hi{क्या} \hi{है}? — \hi{हमारे} \hi{घर} \hi{में} \hi{सब} \hi{है।}
\end{quote}

Six lines, three different owners, and the same bend in every one of them: \emph{āpke}, \emph{tumhāre}, \emph{hamāre} — every possessive ending in \textbf{-\hi{ए}} because a postposition stands after it.
\end{tcolorbox}

\subsection*{Your turn}
Say these aloud:

\begin{itemize}
\raggedright
  \item \emph{āpke pās kyā hai?} and answer it about yourself
  \item the same question to a friend, using the familiar possessive (\emph{tumhāre pās})
  \item \emph{apne dost ke sāth}, and say why it is \emph{apne} rather than \emph{apnā}
  \item one true sentence about something you have
\end{itemize}

\begin{tcolorbox}[breakable,colback=teal!4,colframe=teal!35!black,title={Before you move on}]
Which word in \emph{mere pās ek kitāb hai} means \emph{have}? (None — the construction is \emph{near me there is}.) Why is it \emph{āpke pās} rather than \emph{āpkā pās}? (\emph{Pās} is a postposition, so the possessive goes oblique.)
\end{tcolorbox}
